%%
% Project report – forceps, shadow and rod tracking in a retina simulator
% Template: TUM Corporate Design, class tum-article
% (https://gitlab.lrz.de/tum-templates/templates)
%
% Build:  cd report && latexmk -pdf main.tex
%
% Previous full draft: git show 4f9d46a:report/main.tex
%%

\documentclass[english]{tum-article}
% Variants:
%\documentclass[english, twocolumn]{tum-article}
%\documentclass[english, times]{tum-article}

\usepackage{booktabs}
\usepackage{siunitx}
\usepackage{xcolor}

% Marks open items – remove all before submission
\newcommand{\todo}[1]{\textcolor{TUMOrange}{\textbf{[TODO:} #1\textbf{]}}}

\title{Real-Time Tracking of Micro-Forceps, Shadow and Rod in a Retina Simulator}
\subtitle{Project Report}
\author{Luis Carilla}
\titlerunning{Forceps Tracking in a Retina Simulator}
\authorrunning{L. Carilla}

% \email{...@tum.de}
\affil{\todo{Chair / Department}, Technical University of Munich}
\date{\today}

\begin{document}

\maketitle

\begin{abstract}
  \todo{5--8 sentences: problem, approach, key results. Write last.}
\end{abstract}

% ---------------------------------------------------------------------------
\section{Introduction}
\label{sec:intro}

In retinal surgery, the microscope gives little depth perception. Surgeons judge
the distance to the retina by the instrument's shadow~\todo{cite}.

We track the forceps and its shadow in the video of a retina simulator,
reconstruct the 3D position of the tips and show their distance to the retina in
real time. We also track a rod and detect when the forceps grasps it. The goal is to make such tasks easier to support with human
assistance and, in the future, to automate.

% ---------------------------------------------------------------------------
\section{Background and Related Work}
\label{sec:related}

\todo{Depth from instrument shadows in ophthalmic surgery.}

\todo{Surgical instrument detection and tracking (segmentation, keypoints).}

\todo{Guidance displays / depth feedback for surgeons.}

% ---------------------------------------------------------------------------
\section{Setup and System Overview}
\label{sec:system}

\subsection{Retina Simulator and Footage}
\label{sec:system-sim}

\todo{Simulator, task levels, camera / capture, trajectory logs.}

\subsection{Pipeline Architecture}
\label{sec:system-pipeline}

\todo{Services detector $\rightarrow$ orchestrator $\rightarrow$ stream
  $\rightarrow$ viz; sources (file, RTSP, capture card); data flow and latency.
  Architecture figure.}

% ---------------------------------------------------------------------------
\section{Detecting the Forceps and Its Shadow}
\label{sec:detection}

\todo{Task definition: two jaw tips, two shadow tips, junction.}

\subsection{First Attempt: Segmentation on Synthetic Data}
\label{sec:detection-yolo}

\todo{YOLO segmentation trained on synthetic data (Jul--Aug); what worked,
  why it was abandoned.}

\subsection{Classical Detector}
\label{sec:detection-classical}

\todo{Preprocessing, eye-circle normalisation, background model.}

\todo{Forceps and shadow masks.}

\todo{Geometry: tips and junction.}

\subsection{Learned Keypoint Detector}
\label{sec:detection-learned}

\todo{Model (ResNet34 U-Net heatmaps), labels, training, drop-in backend.}

% ---------------------------------------------------------------------------
\section{3D Reconstruction}
\label{sec:reconstruction}

\subsection{Geometry: Eye, Trocars and Light}
\label{sec:reconstruction-geometry}

\todo{Eye sphere, trocar entry points, light position; tip position from
  instrument and shadow.}

\subsection{From an Orthographic to a Pinhole Camera}
\label{sec:reconstruction-camera}

\todo{Orthographic assumption, its error, pinhole model
  (\texttt{camera\_distance}).}

\subsection{Calibration}
\label{sec:reconstruction-calibration}

\todo{Calibration from the simulator log (trocars, eye radius, jaw length);
  eye tilt between sessions; need for per-session calibration.}

% ---------------------------------------------------------------------------
\section{Rod Tracking and Contact Detection}
\label{sec:rod}

\todo{Rod and rod-shadow detection (classical and learned).}

\todo{Image-based contact / grasp detection, hysteresis against noisy frames.}

\todo{Holding a lost rod, adaptive filter for the resting rod, rod
  calibration with the pinhole camera.}

% ---------------------------------------------------------------------------
\section{Visualization and Guidance}
\label{sec:viz}

\todo{Purpose, users and the questions the workstation answers.}

\todo{Alert logic: state model, thresholds, stabilisation.}

\todo{Dropout and stale-data handling; key design decisions (figures).}

% ---------------------------------------------------------------------------
\section{Evaluation}
\label{sec:evaluation}

\subsection{Ground Truth from Simulator Logs}
\label{sec:eval-gt}

\todo{Trajectory logs as ground truth; alignment of log and video.}

\subsection{2D Detection Accuracy}
\label{sec:eval-2d}

\todo{Classical vs.\ learned: pixel error, runtime, failure cases.}

\subsection{3D Accuracy}
\label{sec:eval-3d}

\todo{Error budget: detection vs.\ camera model; orthographic vs.\ pinhole.}

\subsection{Rod Contact and Alarms}
\label{sec:eval-rod-alarms}

\todo{Contact detection accuracy; alarm replay against ground truth.}

% ---------------------------------------------------------------------------
\section{Discussion and Limitations}
\label{sec:discussion}

\todo{What worked, what did not, limitations of simulator-only evaluation.}

% ---------------------------------------------------------------------------
\section{Conclusion and Outlook}
\label{sec:conclusion}

\todo{Summary and next steps.}

% \section*{Acknowledgements}

The layout follows the TUM corporate design template~\cite{tumtemplates}.

\bibliographystyle{IEEEtran}
\bibliography{literature}

\end{document}
